% !TeX root = ../navier-stokes-manuscript.tex
% ============================================================
% 2.1 Vortex Stretching
% ============================================================

The momentum equation has four pieces:
\[
  \underbrace{\partial_t u}_{\text{local acceleration}}
  +
  \underbrace{(u\cdot\nabla)u}_{\text{nonlinear transport}}
  =
  \underbrace{-\nabla p}_{\text{pressure}}
  +
  \underbrace{\nu\Delta u}_{\text{viscous diffusion}}.
\]
The central regularity problem lies in the competition between nonlinear transport and viscosity.
The hope is simple: viscosity smooths the fluid faster than nonlinear transport can sharpen it.
If that remains true at every scale, the velocity stays smooth.

\subsection{Vortex Stretching}\label{sub:ThePerfectlyStretchingVortex}

Take a small material fluid element carrying vorticity along one axis. Axial strain lengthens this vortex segment, incompressibility draws its rotating fluid inward as the cross-section contracts, and the aligned stretching contribution speeds its interior rotation. Viscosity acts at the same time, so the vorticity grows only when stretching wins the full balance.

Suppose the strain is locally uniform across the element and extends it along the unit axis $e$. If $L(t)$ is its length, write
\[
  S:=\frac12\bigl(\nabla u+(\nabla u)^{\mathsf T}\bigr),
  \qquad
  Se=\sigma(t)e,
  \qquad
  \sigma(t)>0,
  \qquad
  \frac{\dot L(t)}{L(t)}
  =
  e\cdot Se
  =
  \sigma(t).
\]
The positive rate $\sigma(t)$ extends the segment. Its material volume $\mathcal V$ cannot change, so its transverse area must contract. Define $A(t):=\mathcal V/L(t)$ and the effective radius $r_{\mathrm{eff}}(t):=\sqrt{A(t)/\pi}$. Then
\[
  \nabla\cdot u=0,
  \qquad
  \frac{d}{dt}(A L)=0,
  \qquad
  \frac{\dot A}{A}=-\sigma(t),
  \qquad
  \frac{\dot r_{\mathrm{eff}}}{r_{\mathrm{eff}}}
  =
  -\frac{\sigma(t)}{2}.
\]
As $r_{\mathrm{eff}}$ falls, the rotating fluid is concentrated closer to the axis. With $\omega:=\nabla\times u$, the vorticity carried along that axis satisfies
\[
  \frac{D\omega}{Dt}
  =
  S\omega+\nu\Delta\omega,
  \qquad
  \frac{D}{Dt}
  :=
  \partial_t+u\cdot\nabla,
\]
and under the conditional alignment $\omega=|\omega|e$,
\[
  S\omega
  =
  \sigma(t)\omega,
  \qquad
  \frac12\frac{D|\omega|^2}{Dt}
  =
  \sigma(t)|\omega|^2
  +
  \nu\,\omega\cdot\Delta\omega.
\]
The first term records the spin-up from stretching, while the second records the simultaneous viscous effect. The rotation strengthens as the radius contracts only on intervals where their sum remains positive.

Finite kinetic energy does not by itself stop that local concentration. The energy identity proved in \Cref{prop:ClassicalKineticEnergyIdentity} and the antisymmetric part of $\nabla u$ give
\[
  \norm{u(t)}_2
  \leq
  \norm{u_0}_2
  <\infty,
  \qquad
  \left|\frac12\bigl(\nabla u-(\nabla u)^{\mathsf T}\bigr)\right|_{\mathrm F}
  =
  \frac{|\omega|}{\sqrt2}
  \leq
  |\nabla u|_{\mathrm F}.
\]
The total energy can stay finite while the rotational part of the gradient rises inside the narrowing segment. Its width, rather than its total energy, is where the unresolved competition becomes sharp.

\divider{\NavierStokes{} Scaling}

Fix a time $t_0<T_*$ and call the segment's effective radius $\ell:=r_{\mathrm{eff}}(t_0)$. For $\lambda>1$, the exact \NavierStokes{} comparison at the finer width $\lambda^{-1}\ell$ is
\[
  u_\lambda(x,t)
  :=
  \lambda u(\lambda x,\lambda^2t),
  \qquad
  p_\lambda(x,t)
  :=
  \lambda^2p(\lambda x,\lambda^2t).
\]
At time $\lambda^{-2}t_0$, the corresponding segment in $u_\lambda$ has effective radius $\lambda^{-1}\ell$. This rescaled field is another solution at another scale, not a later state of the tracked material segment. Each momentum term acquires the same factor:
\[
\begin{aligned}
  \partial_tu_\lambda(x,t)
  &=
  \lambda^3(\partial_tu)(\lambda x,\lambda^2t),
  \\
  (u_\lambda\cdot\nabla)u_\lambda(x,t)
  &=
  \lambda^3\bigl((u\cdot\nabla)u\bigr)(\lambda x,\lambda^2t),
  \\
  \nabla p_\lambda(x,t)
  &=
  \lambda^3(\nabla p)(\lambda x,\lambda^2t),
  \\
  \nu\Delta u_\lambda(x,t)
  &=
  \lambda^3\nu(\Delta u)(\lambda x,\lambda^2t).
\end{aligned}
\]
The factor $\lambda^3$ strengthens nonlinear transport and viscosity together; pressure and local acceleration keep pace with them. Moving to a finer width gives viscosity no relative advantage.

\divider{Critical Height}

Although the four terms remain matched, the ordinary sizes of the rescaled field move in opposite directions. Let $\Lambda:=(-\Delta)^{1/2}$. Since
\[
  \widehat{u_\lambda}(\xi,t)
  =
  \lambda^{-2}\widehat u(\lambda^{-1}\xi,\lambda^2t),
\]
a change of variables gives
\[
  \norm{u_\lambda(t)}_{\dot H^s}^2
  =
  \lambda^{2s-1}
  \norm{u(\lambda^2t)}_{\dot H^s}^2.
\]
At $s=0$, kinetic energy loses one power of $\lambda$. At $s=1$, the first derivatives gain one. Between them, the exponent vanishes at $s=\tfrac12$, so set
\[
  H(t)
  :=
  \frac12\norm{u(t)}_{\dot H^{1/2}}^2
  =
  \frac12\norm{\Lambda^{1/2}u(t)}_2^2.
\]
This is a global quantity of the velocity field, not a height attached to the vortex core. Writing $H_\lambda$ for the same quantity evaluated on $u_\lambda$,
\[
  \norm{u_\lambda(t)}_2^2
  =
  \lambda^{-1}\norm{u(\lambda^2t)}_2^2,
  \qquad
  H_\lambda(t)=H(\lambda^2t),
  \qquad
  \norm{\nabla u_\lambda(t)}_2^2
  =
  \lambda\norm{\nabla u(\lambda^2t)}_2^2.
\]
Energy falls and the first-derivative norm rises, while $H$ stays unchanged by the spatial rescaling. This half-derivative level is the critical height, and its sign of change along one solution returns us to the nonlinear--viscous balance.

\divider{Nonlinear Production versus Viscous Dissipation}

At critical height, the nonlinear term contributes the signed production
\[
  \mathcal P_{1/2}[u](t)
  :=
  -\operatorname{Re}
  \pair
    {\Lambda^{1/2}u(t)}
    {\Lambda^{1/2}\bigl((u(t)\cdot\nabla)u(t)\bigr)}_{L^2}.
\]
The pressure pairing vanishes because $\nabla\cdot u=0$, while the Laplacian contributes $-\nu\norm{\Lambda^{3/2}u}_2^2$. Hence
\[
  H'(t)
  +
  \nu\norm{\Lambda^{3/2}u(t)}_2^2
  =
  \mathcal P_{1/2}[u](t).
\]
The signed difference is exact: $H$ rises only when nonlinear production exceeds the viscosity dissipating it at the same time. Under the same rescaling,
\[
  \mathcal P_{1/2}[u_\lambda](t)
  =
  \lambda^2\mathcal P_{1/2}[u](\lambda^2t),
  \qquad
  \nu\norm{\Lambda^{3/2}u_\lambda(t)}_2^2
  =
  \lambda^2\nu\norm{\Lambda^{3/2}u(\lambda^2t)}_2^2.
\]
Both rates acquire the same $\lambda^2$ factor. Neither gains at finer scale, and the reciprocal contraction $t\mapsto\lambda^{-2}t$ identifies the amount of comparison time belonging to that change of width.

\divider{The Heat Clock}

For the linear heat equation $v_t=\nu\Delta v$, a Fourier mode evolves by
\[
  \widehat v(\xi,t+\delta t)
  =
  e^{-\nu|\xi|^2\delta t}\widehat v(\xi,t).
\]
A structure of width $r$ is concentrated near frequencies $|\xi|\simeq r^{-1}$. Viscous decay changes those frequencies by order one when
\[
  \nu|\xi|^2\delta t\simeq1,
\]
which gives the linear comparison time
\[
  \tau_\nu(r)
  :=
  \frac{r^2}{\nu}.
\]
This is a heat-flow clock for a width, not a proved lifetime of the tracked vortex. Halving $r$ quarters the clock, and more generally
\[
  \tau_\nu(\lambda^{-1}r)
  =
  \lambda^{-2}\tau_\nu(r).
\]
The same inverse-square contraction that scaled the full equation now attaches a shorter comparison interval to every finer width.

\divider{The Zeno Cascade}

Repeat the arithmetic halving. For
\[
  r_n:=2^{-n}r_0,
  \qquad
  \tau_n:=\frac{r_n^2}{\nu},
\]
the corresponding heat clocks satisfy
\[
  \tau_n
  =
  \frac{r_0^2}{\nu}4^{-n},
  \qquad
  \sum_{n=0}^{\infty}\tau_n
  =
  \frac{4r_0^2}{3\nu}
  <
  \infty.
\]
The convergent sum is only arithmetic. To describe a candidate singular history, one \NavierStokes{} solution must carry the original tracked material segment through widths
\[
  r_0>r_1>r_2>\cdots\longrightarrow0
\]
at times
\[
  t_0<t_1<t_2<\cdots\uparrow T_*<\infty,
  \qquad
  t_{n+1}-t_n\asymp\tau_n,
\]
with $r_{\mathrm{eff}}(t_n)\asymp r_n$ and all comparison constants independent of $n$. Only then does the remaining-time calculation belong to one material history:
\[
  T_*-t_n
  \asymp
  \sum_{k=n}^{\infty}\tau_k
  =
  \frac{4r_n^2}{3\nu}.
\]
At the $n$-th narrowing, both the next transition and the total time left are of order $r_n^2/\nu$. One velocity field would have to produce each successive contraction of the tracked segment on intervals that collapse to zero as $t\uparrow T_*$.
